% SPDX-License-Identifier: CC-BY-SA-4.0
% Part: Lexical Analysis
% Section: Arden's lemma and identities
% Exercise: Solution of a system of regular equations

\begingroup

\begin{exercice}[Systèmes d'équations rationnelles]\label{exo:lexing/arden/arden}

  Résoudre les systèmes d'équations rationnelles suivants :
  \begin{question}
  \item $\left\{\begin{array}{rcl}
    L_1 &=& aL_2 \mid  bL_3 \\
    L_2 &=& bL_2 \mid  \varepsilon \\
    L_3 &=& bL_3 \mid  aL_2
  \end{array}\right.$
  \item $\left\{\begin{array}{rcl}
    L_1 &=& bL_2 \mid  bL_3 \\
    L_2 &=& aL_2 \mid  aL_4 \mid  \varepsilon \\
    L_3 &=& cL_4 \\
    L_4 &=& cL_4 \mid  \varepsilon
  \end{array}\right.$
  \end{question}

\end{exercice}

\begin{correction}

  \begin{question}
  \item
    $\varepsilon b\notin b$, donc par le lemme d'Arden, l'équation
    $L_2 = b L_2 | \varepsilon$ a une unique solution $L_2 = b^\star \varepsilon = b^\star$\\
    $L_3 = bL_3 \mid  ab^\star  = b^\star  ab^\star $ (par le lemme d'arden)\\
    $L_1 = aL_2 \mid  bL_3 = ab^\star  \mid  bb^\star  ab^\star  = (\varepsilon |bb^\star ) ab^\star = b^\star  ab^\star $
  \item
    $L_4 = c^\star $ (par Arden)\\
    $L_3 = c^+$ (par Arden)\\
    $L_2 = a^\star (ac^\star  \mid  \varepsilon)$ (par Arden) $= a^+c^\star  \mid  a^\star  = a^+c^\star  \mid  a^+ \mid  \varepsilon = a^+c^\star  \mid  \varepsilon$\\
    $L_1 = b(a^+c^\star  \mid  \varepsilon) \mid  bc^+ = b(a^+c^\star  \mid  \varepsilon \mid  c^+) = b(a^+c^\star  \mid  c^\star ) = ba^\star  c^\star $
  \end{question}

\end{correction}

\endgroup
\endinput
